%......................................................................
\subsection{Logistic regression as a single neuron}
\label{sec:logregneuron}

\index{perceptron}\index{artificial neuron}\index{delta rule}
Before we take the class apart with automatic differentiation, it is worth
drawing it.  Written out for one observation with $p$ features, the model
of Eq.~(\ref{eq:5-logodds}) is
\begin{equation}
  z = \theta_0 + \sum_{j=1}^{p}\theta_j x_j ,
  \qquad
  \hat p = \sigma(z) = P(y=1\mid\bm{x}) ,
  \label{eq:5-neuron}
\end{equation}
a weighted sum of the inputs, shifted by the intercept and passed through
one non-linear function.  Figure~\ref{fig:5-neuron} shows
Eq.~(\ref{eq:5-neuron}) as a diagram: the features enter on the left, each
edge carries a parameter, the node $\Sigma$ forms the linear predictor $z$,
the node $\sigma$ squashes it into a probability, and, if we also attach
the label, the cross entropy of Eq.~(\ref{eq:5-crossentropycompact}) is one
more node at the far right.  This diagram is the \emph{simple perceptron}
of Rosenblatt, mentioned in Section~\ref{sec:classification}, with one
change: the step function that made the perceptron a hard classifier has
been replaced by the sigmoid.  In the language of Chapter~\ref{chap:nn},
which begins with this very picture, the $x_j$ are the inputs of an
\emph{artificial neuron}, the $\theta_j$ are its \emph{weights}, the
intercept $\theta_0$ is its \emph{bias}, $\sigma$ is its \emph{activation
function} and $z$ its \emph{pre-activation}.  Logistic regression is a
neural network with a single neuron and no hidden layer; nothing in this
chapter needs to be renamed when we get there.

\begin{figure}[htbp]
\centering
\begin{tikzpicture}[>=Latex,
  inp/.style={circle,draw=black!65,fill=Blue3!12,minimum size=8.5mm,
              inner sep=0pt,font=\small},
  bia/.style={circle,draw=black!65,fill=Red3!12,minimum size=8.5mm,
              inner sep=0pt,font=\small},
  sm/.style={circle,draw=black!65,fill=black!5,minimum size=10mm,inner sep=0pt},
  ac/.style={rectangle,rounded corners=1pt,draw=black!65,fill=green!12,
             minimum height=10mm,minimum width=10mm,inner sep=1pt},
  ls/.style={rectangle,rounded corners=1pt,draw=black!65,dashed,fill=black!3,
             minimum height=10mm,minimum width=13mm,inner sep=2pt,font=\small},
  wl/.style={font=\footnotesize,fill=white,inner sep=1pt}]
  \node[inp] (x1) at (0, 2.10) {$x_1$};
  \node[inp] (x2) at (0, 1.10) {$x_2$};
  \node      (xd) at (0, 0.30) {$\vdots$};
  \node[inp] (xp) at (0,-0.55) {$x_p$};
  \node[sm] (S) at (3.60,0.78) {$\sum$};
  \node[ac] (F) at (5.60,0.78) {$\sigma$};
  \node[bia] (b) at (3.60,-1.35) {$1$};
  \node[ls] (L) at (8.40,0.78) {$\ln(1+e^{z})-yz$};
  \node[inp] (y) at (8.40,-1.35) {$y$};
  \draw[->,black!70] (x1) -- node[wl,pos=0.62]{$\theta_1$} (S);
  \draw[->,black!70] (x2) -- node[wl,pos=0.62]{$\theta_2$} (S);
  \draw[->,black!70] (xp) -- node[wl,pos=0.62]{$\theta_p$} (S);
  \draw[->,black!70] (b)  -- node[wl,pos=0.5]{$\theta_0$} (S);
  \draw[->,black!70] (S) -- node[above,font=\footnotesize]{$z$} (F);
  \draw[->,black!70] (F) -- node[above,font=\footnotesize]{$\hat p=\sigma(z)$} (L);
  \draw[->,black!70] (y) -- (L);
  \draw[->,black!70] (L) -- ++(1.6,0) node[right,font=\small] {$C_i$};
  \node[font=\footnotesize,align=center] at (3.60,2.35)
       {$z=\theta_0+\displaystyle\sum_{j=1}^{p}\theta_jx_j$};
  \node[font=\footnotesize,black!60] at (0,-1.35) {inputs};
  \node[font=\footnotesize,black!60] at (5.60,-1.35) {activation};
  \node[font=\footnotesize,black!60] at (8.40,2.35) {loss};
  \draw[<-,black!50,dashed] ([yshift=-3pt]F.south) -- ++(0,-0.45)
       node[below,font=\footnotesize,black!60,align=center]
       {reverse sweep:\\ $\bar z=\hat p-y$};
\end{tikzpicture}
\caption{Logistic regression drawn as a single artificial neuron,
Eq.~(\ref{eq:5-neuron}): the features $x_j$ enter on the left, the edges
carry the parameters $\theta_j$, the intercept $\theta_0$ is the weight of
an input permanently fixed at one, the node $\Sigma$ forms the linear
predictor $z$ and the activation $\sigma$ turns it into the probability
$\hat p$.  Attaching the label $y$ and the cross entropy of
Eq.~(\ref{eq:5-crossentropycompact}) on the right gives the cost of one
observation.  With the step function in place of $\sigma$ this is
Rosenblatt's perceptron.  Read from right to left, the same diagram is the
reverse sweep of Section~\ref{sec:logregad}: the sensitivity of the cost to
the pre-activation is the residual $\hat p-y$, and the gradient with
respect to $\theta_j$ is that residual times the input $x_j$ on the
corresponding edge.}
\label{fig:5-neuron}
\end{figure}

\paragraph{The perceptron learning rule.}
The picture makes the gradient obvious.  The cost of one observation
depends on the parameters only through $z$, so by the chain rule
$\partial C_i/\partial\theta_j
 =(\partial C_i/\partial z)(\partial z/\partial\theta_j)$.  The first factor
is the derivative of $\ln(1+e^{z})-y_iz$, namely $\sigma(z)-y_i=\hat p_i-y_i$,
the summand of Eq.~(\ref{eq:5-dtheta0}); the second is the input $x_{ij}$
sitting on the edge of $\theta_j$ (with $x_{i0}=1$ for the bias).  Hence
\begin{equation}
  \frac{\partial C_i}{\partial\theta_j} = (\hat p_i-y_i)\,x_{ij},
  \qquad
  \nabla_{\bm{\theta}}C_i = (\hat p_i-y_i)\,\bm{x}_i ,
  \label{eq:5-gradone}
\end{equation}
and a gradient step on this one observation reads
\begin{equation}
  \theta_j \;\leftarrow\; \theta_j - \gamma\,(\hat p_i-y_i)\,x_{ij} .
  \label{eq:5-deltarule}
\end{equation}
Equation~(\ref{eq:5-deltarule}) is, up to the choice of activation,
Rosenblatt's original perceptron learning rule of 1958, and with a
differentiable activation it is the \emph{delta rule} of Widrow and Hoff:
increase the weight on every input that was active when the output was too
small, decrease it when the output was too large, in proportion to the
error.  Summing Eq.~(\ref{eq:5-gradone}) over the observations gives back
$\bm{X}^{T}(\bm{p}-\bm{y})$, the gradient~(\ref{eq:5-gradient}) we
derived by algebra in Section~\ref{sec:logreggradient}.  Applying
Eq.~(\ref{eq:5-deltarule}) one observation, or a few observations, at a
time rather than summing first is stochastic gradient descent, to which
Section~\ref{sec:logregsgd} is devoted.

\paragraph{From the neuron to the computational graph.}
Figure~\ref{fig:5-neuron} can be read in two ways, and the difference
between them is the subject of the next subsection.  Read as a
\emph{model}, information flows from left to right: the data $\bm{x}$ enter,
the parameters are fixed numbers on the edges, and a probability comes out.
Read as a \emph{computational graph} in the sense of
Section~\ref{sec:adtrace}, it is the same nodes and the same edges, but the
roles are exchanged: the parameters are the inputs with respect to which we
differentiate, the data are constants of the trace, and the output is the
scalar cost $C_i$.  Automatic differentiation in reverse mode is then
nothing but the neuron read backwards.  The adjoint of the cost is one; it
passes through the loss node to give the sensitivity of $C_i$ to $\hat p$,
through the activation to give the sensitivity to the pre-activation,
$\bar z=\hat p-y$, and from there along each edge to the parameter it
carries, picking up the factor $x_j$ on the way -- which is
Eq.~(\ref{eq:5-gradone}) once more, now produced by the fan-out rule rather
than by the chain rule written out by hand.  The next subsection performs
exactly this sweep on the trace of the graph, node by node, and compares it
with the forward sweep.  The reason for insisting on the picture is what
comes after: in Chapter~\ref{chap:nn} the inputs $x_j$ of this neuron will
themselves be the outputs of other neurons, the reverse sweep will not stop
at $\bar z$ but continue along the edges into the layer before, and
Eq.~(\ref{eq:5-gradone}) will be applied layer by layer.  That is
backpropagation, and Figure~\ref{fig:5-neuron} is its one-neuron case.

\begin{notebox}
\textbf{Why the sigmoid, and not the step.}  The perceptron and logistic
regression share Figure~\ref{fig:5-neuron}, and the reason the latter
replaced the former is visible in it.  With the step function as
activation the node $\sigma$ has derivative zero almost everywhere, the
reverse sweep dies at that node, and no gradient reaches the parameters:
Rosenblatt's rule had to be justified by a separate convergence argument
that works only for separable data.  With the sigmoid every node of the
graph is differentiable, the sweep passes through, and the rule becomes an
instance of gradient descent on a convex cost, with all of
Chapter~\ref{chap:optim} behind it.  The same requirement -- that the
activation be differentiable so that the sweep can pass through it -- is
what shapes the choice of activation functions in Chapter~\ref{chap:nn}.
\end{notebox}

